\subsection{教育项目推荐与Gen-AI整合策略}
\label{subsec:task3_programs}

基于前述职业选择及其任务层面的Gen-AI影响分析，我们为三个代表性职业推荐具体的教育项目。
推荐遵循三重准则：\textbf{（1）机构层面的行业领导力}、\textbf{（2）课程与产业需求的对齐度}、以及\textbf{（3）Gen-AI资源整合的制度化路径}。
具体推荐项目及其详细选择依据见附录表~\ref{tab:programs_genai_rationale}。

针对\textbf{机器人工程师}，我们推荐卡内基梅隆大学机器人研究院的MRSD项目，该项目具备行业领导力且已建立Gen-AI课程资源接口。
针对\textbf{机械维修一线主管}，Universal Technical Institute（UTI）展现出快速响应产业技术迁移的能力，并已将AI诊断纳入教学叙事。
针对\textbf{法庭速录师}，Generations College作为美国首创法庭速录项目的历史传承者，其Realtime训练体系可与Gen-AI后编辑工具形成互补。

\subsubsection{基于文本的课程剪枝与增补元模型}

在完成项目选择后，我们进一步构建一个\textbf{纯文本驱动的课程优化元模型}，用于指导各教育项目如何在Gen-AI时代进行课程结构调整。该模型仅依赖课程标题与描述文本，无需学分信息，具备高度可解释性与可迁移性。

\paragraph{三轴评估框架}
我们为每门课程 $m$ 构造三个可解释的代理指标：

\textbf{（1）就业轴 $E_m$}：衡量课程文本与劳动力市场技能需求的对齐程度。该指标通过统计课程描述中命中预定义就业关键词词典 $\mathcal{K}_E$ 的频次并归一化获得：
\\这部分是必要的，不能过于理想主义，仍然需要结合真实的社会形态。
\begin{equation}
E_m = \text{norm}\left(\sum_{w \in \mathcal{K}_E} \mathbf{1}[w \in t_m]\right)
\end{equation}
其中 $t_m$ 为课程 $m$ 的文本内容（标题 + 描述）。该指标的作用类似于\textbf{职业相关性过滤器}，识别与产业技能直接挂钩的课程模块。

\textbf{（2）人文轴 $H_m$}：衡量课程对批判性思维、伦理意识、跨学科素养等人文能力的贡献程度：
\begin{equation}
H_m = \text{norm}\left(\sum_{w \in \mathcal{K}_H} \mathbf{1}[w \in t_m]\right)
\end{equation}
人文轴的引入是对"唯就业论"的关键修正——我们认为，Gen-AI时代恰恰需要强化人类独有的判断力与价值观锚定能力。

\textbf{（3）冗余度 $R_m$}：捕捉课程与同一项目中其他课程的内容重叠程度：
\begin{equation}
R_m = \max_{n \neq m} \text{sim}(t_m, t_n)
\end{equation}
该指标识别"信息熵最低"的模块——即那些可被其他课程替代的冗余内容。

\paragraph{剪枝决策规则}
我们采用\textbf{三重门槛联合约束}的保守剪枝策略：仅当一门课程同时满足"就业相关性低"、"人文贡献低"、"与其他课程高度冗余"三个条件时，才建议将其从课程体系中移除：
\begin{equation}
\text{PRUNE}(m) = \mathbf{1}\left[(E_m < e_{\min}) \land (H_m < h_{\min}) \land (R_m > r_{\min})\right]
\end{equation}
课程保留决策变量定义为 $x_m = 1 - \text{PRUNE}(m)$。这种设计确保了我们不会误删那些"就业导向弱但人文价值高"的课程——例如伦理学、科技史等在AI时代反而更加重要的模块。

\paragraph{增补建议机制}
在剪枝的同时，模型检测整体人文轴均值是否低于目标区间：
\begin{equation}
\overline{H} = \frac{1}{M} \sum_{m} H_m
\end{equation}
若 $\overline{H}$ 低于阈值，模型建议新增特定类型的模块集合 $\Delta\mathcal{M}$，以提升课程体系的人文素养覆盖度。增补建议的具体内容依据各职业的Gen-AI影响特征动态生成。

\paragraph{模型局限与缓解措施}
当课程描述过短或偏向营销性语言时，$E_m$ 与 $H_m$ 可能被低估。我们通过以下方式缓解该问题：（1）允许手动扩充关键词词典；（2）报告阈值参数的敏感性区间，供决策者参考调整。

